\subsection{Full-Stack TGN Detector: Offline Validation Results (Full-Scale, 300-Epoch Dataset)}
\label{sec:tgn-offline-validation-v2}

This subsection reports the offline train/validation/test evaluation of the
Full-Stack TGN detector (\S\ref{sec:tgn-arch}, Eq.~3.25--3.27) on the
current, corrected training corpus, superseding the earlier pilot result
reported in \S\ref{sec:tgn-offline-validation}. The corpus was regenerated
after correcting the sliding-window bound derivation (edge-freshness decay
rate $\gamma$ and $L_{\text{link}}$, Eq.~3.22), and the reported model was
retrained end-to-end against the corrected values.

\paragraph{Compute environment.}
Training was performed with PyTorch (CUDA build) on an NVIDIA GPU. The
reported run used up to 3 training attempts (restart target val~MCC~$\geq
0.975$) at 250 epochs per attempt; the target was met on attempt 3. Total
wall-clock time for the reported run was 12{,}803.0\,s ($\approx$213.4\,min,
$\approx$3.56\,GPU-hours).

\paragraph{Dataset generation.}
Events were regenerated from all twelve attack scenarios via the NS-3.35
SDVN simulation (\texttt{scratch/routing.cc}), with
$\texttt{sim\_time}=60\,\text{s}$, $N_{\text{vehicles}}=200$,
$N_{\text{controllers}}=4$, $N_{\text{RSUs}}=64$ for RSU-bearing scenarios
(0 otherwise), simulation seed $\texttt{RngRun}=1$ (independent of the
model seed below), and an $\texttt{attack\_percentage}$ sweep over
$\{0, 20, 40, 60, 80, 100\}\%$ per scenario, combined into a single corpus.
A total of 42{,}151 non-BEACON events were retained (31{,}954 attack,
10{,}197 benign). Table~\ref{tab:tgn-val2-dataset} gives the per-scenario
composition.

\begin{table}[htbp]
\centering
\caption{TGN training dataset: per-scenario event composition.}
\label{tab:tgn-val2-dataset}
\begin{tabular}{@{}llrrr@{}}
\toprule
\# & Scenario & Events & Attack & Benign \\
\midrule
1  & TTW-S1  & 13{,}652 & 11{,}178 & 2{,}474 \\
2  & TTW-S2  & 6{,}134  & 4{,}267  & 1{,}867 \\
3  & TTW-S3  & 890      & 180      & 710     \\
4  & TTW-S4  & 916      & 182      & 734     \\
5  & BSHH-S1 & 7{,}516  & 6{,}345  & 1{,}171 \\
6  & BSHH-S2 & 10{,}780 & 8{,}717  & 2{,}063 \\
7  & BSHH-S3 & 709      & 351      & 358     \\
8  & BSHH-S4 & 631      & 315      & 316     \\
9  & ME-S1   & 12       & 0        & 12      \\
10 & ME-S2   & 467      & 0        & 467     \\
11 & ME-S3   & 221      & 212      & 9       \\
12 & ME-S4   & 223      & 207      & 16      \\
\midrule
\multicolumn{2}{l}{\textbf{Total}} & \textbf{42{,}151} & \textbf{31{,}954} & \textbf{10{,}197} \\
\bottomrule
\end{tabular}
\end{table}

\paragraph{Data split.}
The same stratified, per-(scenario, class) temporal split procedure as the
pilot run (\S\ref{sec:tgn-pipeline}) was applied. Table~\ref{tab:tgn-val2-split}
gives the resulting fold sizes.

\begin{table}[htbp]
\centering
\caption{TGN train/validation/test split sizes.}
\label{tab:tgn-val2-split}
\begin{tabular}{@{}lrrr@{}}
\toprule
Split & Events & Attack & Benign \\
\midrule
Train      & 29{,}492 & 22{,}361 & 7{,}131 \\
Validation & 6{,}322  & 4{,}794  & 1{,}528 \\
Test       & 6{,}337  & 4{,}799  & 1{,}538 \\
\bottomrule
\end{tabular}
\end{table}

\paragraph{Preprocessing and features.}
The five-element feature vector
$x_v = [\tau_{\text{dev}}, c_v^{W}, \Delta s_v, \rho_v, \iota_v] \in \mathbb{R}^5$
(Eq.~3.20), plus the separate time-elapsed encoding $\phi$ (Eq.~3.21),
beacon interval $T_b=0.1$\,s. The edge-freshness decay rate $\gamma$ and
sliding-window bound $W_{\max}$ are auto-derived from the corrected
$L_{\text{link}}=20.4$\,s (Eq.~3.22), giving $\gamma=147.2$ and
$W_{\max}=204$ for this run — superseding the earlier
$L_{\text{link}}=43.0$\,s/$\gamma=310.2$/$W_{\max}=430$ values used in the
pilot, which were derived from a stale intermediate value.

\paragraph{Model architecture.}
GRU-based per-node memory with embedding dimension $d=192$ and $L=2$
rounds of message passing over the 3-node induced subgraph
$\{\text{reporter}, \text{link\_src}, \text{link\_dst}\}$ (the GRU cell
itself is a single, non-stacked cell; $L$ governs message-passing rounds,
not GRU depth). The GRU's actual input dimension is the concatenation of
the hidden state, the 5-element feature vector $x_v$ (Eq.~3.20), and the
separate $\phi$ term (Eq.~3.21):
$h_{\text{GRU}} = d + 5 + 1 = 192 + 5 + 1 = 198$. Per-node truncated
backpropagation-through-time window $W_{\text{BPTT}}=50$ events (selected
by the sensitivity check below, superseding an earlier
$W_{\text{BPTT}}=100$ configuration). Joint output heads: binary
attack-score logit and 3-way attack-family classifier (TTW/BSHH/ME).

\paragraph{Training hyperparameters.}
Table~\ref{tab:tgn-val2-hparams} lists the values used for the reported run.

\begin{table}[htbp]
\centering
\caption{Hyperparameters used for the reported TGN run.}
\label{tab:tgn-val2-hparams}
\begin{tabular}{@{}llp{7cm}@{}}
\toprule
Parameter & Value & Notes \\
\midrule
Embedding dim $d$ & 192 & \\
GRU input dim $h_{\text{GRU}}$ & 198 & $d + 5 + 1$ (hidden state + Eq.~3.20 5-element feature vector + Eq.~3.21 $\phi$ term); GRU cell is single, non-stacked \\
Message-passing rounds $L$ & 2 & over 3-node induced subgraph \\
Per-node TBPTT window $W_{\text{BPTT}}$ & 50 events & selected over 100 (earlier config) and 200 by the sensitivity check below \\
Epochs & 250 per attempt & up to 3 attempts \\
Optimiser & Adam & learning rate $1\times10^{-3}$ \\
Loss: binary term & BCEWithLogitsLoss & \texttt{pos\_weight} override = 0.3191 \\
Loss: variant CE term & weight 0.3 & 3-way TTW/BSHH/ME classifier \\
Seed & 5 & fixed \\
Restart target & val MCC $\geq 0.975$ & met on attempt 3/3 \\
Winning run & attempt 3/3, epoch 250 of 250 & val\_bestMCC=0.977, val\_AUROC=0.999 \\
Threshold $\theta_{FS}$ & 0.193 & plateau-midpoint MCC maximisation on validation set \\
\bottomrule
\end{tabular}
\end{table}

\paragraph{Sensitivity check: message-passing rounds $L$ and TBPTT window
$W_{\text{BPTT}}$.} Each axis was independently varied away from an
initial baseline configuration ($d=192$, \texttt{pos\_weight}=0.3191,
\texttt{ce\_weight}=0.3, seed=5, $L=2$, $W_{\text{BPTT}}=100$; 250 epochs
per attempt, up to 3 restarts, on the same corrected dataset). The $L$
axis ($L\in\{1,2,3\}$) was swept first, holding $W_{\text{BPTT}}=100$
fixed; the $W_{\text{BPTT}}$ axis ($W_{\text{BPTT}}\in\{50,100,200\}$) was
then swept holding $L=2$ fixed. Table~\ref{tab:tgn-val2-lwbptt-sweep}
reports both, with the selected configuration marked.

\begin{table}[htbp]
\centering
\caption{Message-passing rounds $L$ and TBPTT window $W_{\text{BPTT}}$ sensitivity check, test MCC.}
\label{tab:tgn-val2-lwbptt-sweep}
\begin{tabular}{@{}llr@{}}
\toprule
Axis (fixed at $W_{\text{BPTT}}=100$) & Value & Test MCC \\
\midrule
$L$ & 1              & 0.749 \\
$L$ & \textbf{2 (selected)} & \textbf{0.958} \\
$L$ & 3              & 0.772 \\
\midrule
Axis (fixed at $L=2$) & Value & Test MCC \\
\midrule
$W_{\text{BPTT}}$ & \textbf{50 (selected)} & \textbf{0.965} \\
$W_{\text{BPTT}}$ & 100 (earlier config) & 0.958 \\
$W_{\text{BPTT}}$ & 200            & 0.952 \\
\bottomrule
\end{tabular}
\end{table}

On the $L$ axis, $L=2$ is confirmed by a wide margin — both alternatives
collapse (test MCC 0.749 for $L=1$, 0.772 for $L=3$, vs.\ 0.958 for
$L=2$), not a marginal call; this sweep was conducted at
$W_{\text{BPTT}}=100$, before the $W_{\text{BPTT}}$ axis below was
re-optimised, but the margin is wide enough (and consistent with an
independent catastrophic-collapse finding for $L=1$/$L=3$ on a precursor
dataset) that it is very unlikely to flip at $W_{\text{BPTT}}=50$; it was
not separately re-verified there. On the $W_{\text{BPTT}}$ axis,
$W_{\text{BPTT}}=50$ (test MCC=0.965) outperforms both $W_{\text{BPTT}}=100$
(0.958) and $W_{\text{BPTT}}=200$ (0.952) — a mild downward trend as the
window widens, with 50 as the clear best of the three points tested. Full
per-variant/per-origin comparison (Table~\ref{tab:tgn-val2-wbptt-compare})
shows this is not just an aggregate artefact: $W_{\text{BPTT}}=50$
improves the previously-weakest TTW/control-plane cell from MCC=0.816 to
0.863, the largest single-cell gain in the comparison.
\textbf{$W_{\text{BPTT}}=50$ is therefore adopted as the reported
configuration}, superseding the earlier $W_{\text{BPTT}}=100$ run; all
results below reflect this model.

\begin{table}[htbp]
\centering
\caption{$W_{\text{BPTT}}=50$ (selected) vs.\ $W_{\text{BPTT}}=100$ (earlier config), full breakdown, test MCC.}
\label{tab:tgn-val2-wbptt-compare}
\begin{tabular}{@{}lrr@{}}
\toprule
Metric & $W_{\text{BPTT}}$=100 & $W_{\text{BPTT}}$=50 (selected) \\
\midrule
Overall             & 0.958 & \textbf{0.965} \\
TTW                  & 0.943 & \textbf{0.949} \\
BSHH                 & 0.978 & \textbf{0.986} \\
ME                   & \textbf{0.958} & 0.958 \\
Data-plane           & 0.957 & \textbf{0.964} \\
Control-plane        & 0.940 & \textbf{0.951} \\
TTW / data-plane     & 0.943 & \textbf{0.948} \\
TTW / control-plane  & 0.816 & \textbf{0.863} \\
BSHH / data-plane    & 0.974 & \textbf{0.984} \\
BSHH / control-plane & 1.000 & 1.000 \\
ME / control-plane   & \textbf{1.000} & 0.888 \\
Variant classifier F1 & 0.725 & \textbf{0.736} \\
\bottomrule
\end{tabular}
\end{table}

$W_{\text{BPTT}}=50$ wins or ties on every metric except ME/control-plane
(1.000$\to$0.888), a small-sample cell ($n_{\text{attack}}=69$) likely
reflecting sample sensitivity rather than a genuine regression, given the
consistent improvement everywhere else.

\paragraph{Assessed results --- overall binary detection.}
Table~\ref{tab:tgn-val2-results} reports the assessed test metrics at
$\theta_{FS}=0.193$, the held-out estimate (test set not used for model
or threshold selection).

\begin{table}[htbp]
\centering
\caption{TGN detector: assessed test metrics, overall (binary attack/benign).}
\label{tab:tgn-val2-results}
\begin{tabular}{@{}lrrrrrrr@{}}
\toprule
Split & MCC & AUROC & Accuracy & TP & TN & FP & FN \\
\midrule
Validation (best, epoch 250) & 0.977 & 0.999 & --    & --    & --    & --  & -- \\
Test (held-out)               & 0.965 & 0.997 & 0.987 & 4{,}751 & 1{,}504 & 34 & 48 \\
\bottomrule
\end{tabular}
\end{table}

\paragraph{Assessed results --- per-variant detection MCC.}
Table~\ref{tab:tgn-val2-variant-mcc} reports family-vs-benign detection MCC
computed on independent test subsets, one per attack family.

\begin{table}[htbp]
\centering
\caption{TGN detector: per-variant (family-vs-benign) test MCC, independent subsets.}
\label{tab:tgn-val2-variant-mcc}
\begin{tabular}{@{}lrrrrrrrr@{}}
\toprule
Family & MCC & AUROC & Accuracy & TP & TN & FP & FN & $n_{\text{attack}}$ \\
\midrule
TTW  & 0.949 & 0.995 & 0.980 & 2{,}326 & 854 & 18 & 48 & 3{,}246 \\
BSHH & 0.986 & 1.000 & 0.996 & 2{,}361 & 575 & 13 & 0  & 2{,}949 \\
ME   & 0.958 & 0.994 & 0.979 & 64      & 75  & 3  & 0  & 142     \\
\bottomrule
\end{tabular}
\end{table}

\paragraph{Assessed results --- per-origin detection MCC.}
Table~\ref{tab:tgn-val2-origin-mcc} splits the test set by attacker
placement: data-plane origin (vehicle/RSU-injected, scenarios S1/S2 of
each family) versus control-plane origin (controller-internal, scenarios
S3/S4 of each family).

\begin{table}[htbp]
\centering
\caption{TGN detector: per-origin (data-plane vs.\ control-plane) test MCC, independent subsets.}
\label{tab:tgn-val2-origin-mcc}
\begin{tabular}{@{}lrrrrrrrr@{}}
\toprule
Origin & MCC & AUROC & Accuracy & TP & TN & FP & FN & $n_{\text{attack}}$ \\
\midrule
Data-plane (S1/S2)    & 0.964 & 0.998 & 0.988 & 4{,}542 & 1{,}179 & 33 & 36 & 5{,}790 \\
Control-plane (S3/S4) & 0.951 & 0.984 & 0.976 & 209     & 325     & 1  & 12 & 547     \\
\bottomrule
\end{tabular}
\end{table}

\paragraph{Assessed results --- cross (per-variant $\times$ per-origin) detection MCC.}
Table~\ref{tab:tgn-val2-cross-mcc} gives the finest-grained breakdown,
each of the six (family, origin) cells computed on an independent test
subset.

\begin{table}[htbp]
\centering
\caption{TGN detector: per-variant $\times$ per-origin test MCC, independent subsets.}
\label{tab:tgn-val2-cross-mcc}
\begin{tabular}{@{}llrrrrrrr@{}}
\toprule
Family & Origin & MCC & AUROC & Accuracy & TP & TN & FP & FN \\
\midrule
TTW  & Data-plane    & 0.948 & 0.996 & 0.982 & 2{,}282 & 635 & 18 & 36 \\
TTW  & Control-plane & 0.863 & 0.945 & 0.956 & 44      & 219 & 0  & 12 \\
BSHH & Data-plane    & 0.984 & 1.000 & 0.995 & 2{,}260 & 473 & 13 & 0  \\
BSHH & Control-plane & 1.000 & 1.000 & 1.000 & 101     & 102 & 0  & 0  \\
ME   & Data-plane    & 1.000$^\dagger$ & 1.000$^\dagger$ & 0.973 & 0       & 71  & 2  & 0  \\
ME   & Control-plane & 0.888 & 1.000 & 0.986 & 64      & 4   & 1  & 0  \\
\bottomrule
\end{tabular}
\\[2pt]
\footnotesize{$^\dagger$This subset contains zero attack-labelled events
(TP=FN=0), making the raw MCC formula's denominator zero (undefined).
Following the $\epsilon$-smoothing convention
$\text{MCC}=(TP{\cdot}TN-FP{\cdot}FN+\epsilon)/(\sqrt{(TP{+}FP)(TP{+}FN)(TN{+}FP)(TN{+}FN)}+\epsilon)$,
this evaluates to 1 as $\epsilon\to0^+$ since both the numerator and
denominator are exactly 0 before smoothing; AUROC is reported by the same
convention. This value does not reflect the 2 false positives recorded
among the 71 true negatives in this subset (Accuracy=0.973, not 1.000).}
\end{table}

The TTW/control-plane cell (MCC=0.863) is the weakest of the six and is
discussed in \S\ref{sec:tgn-limitations} as a structural difficulty
specific to controller-internal TTW: unlike data-plane variants, no
externally observable network-layer packet is injected, giving the model
strictly less direct evidence than every other cell in the table.

\paragraph{Assessed results --- variant classification.}
Table~\ref{tab:tgn-val2-variant} reports the 3-way TTW/BSHH/ME classifier
head (Eq.~3.26) on the test split, restricted to attack-labelled rows.

\begin{table}[htbp]
\centering
\caption{TGN variant classifier: assessed test-split metrics.}
\label{tab:tgn-val2-variant}
\begin{tabular}{@{}lr@{}}
\toprule
Metric & Value \\
\midrule
Accuracy       & 0.799 \\
F1 (macro)     & 0.736 \\
$n$ (rows evaluated) & 6{,}337 \\
\bottomrule
\end{tabular}
\end{table}

\paragraph{Known limitation: ME-S1/ME-S2 still contribute no attack-labelled events.}
Table~\ref{tab:tgn-val2-dataset} confirms this pipeline-architectural
property persists at the current dataset scale: ME-S1 and ME-S2 contribute
0 attack-labelled events out of 12 and 0 out of 467 events respectively.
Tracing the exact interception mechanism: both scenarios' forged echo
reports carry the attacking node's own genuine identity and signature (V3/V4
for ME-S1; the malicious RSU for ME-S2), so the HMAC check (Eq.~3.15)
passes without issue — that is \emph{not} what intercepts them. The actual
gate is the location-binding check (Eq.~3.29, \texttt{TetaGuardLocBindVerify}):
each reporter's own real, signed position is checked against the link it
claims to have observed, and must fall within $R_{\text{comm}}$. An
echo/false-witness report is by definition a node claiming to have
observed a link it was not positioned to observe, so this check fails for
essentially every ME-S1/S2 instance independent of any other factor; a
second gate (Eq.~3.32, witness quorum — a majority of reporters must agree
on the link) would additionally need to be cleared even if location-binding
were somehow satisfied. When Stage~0 rejects an event, the event-emission
path returns immediately and the Stage-1 signature evaluator that computes
the features/labels TGN trains on is never invoked, so the event never
becomes a TGN training row (it may still appear in the raw audit log for
diagnostic purposes). Controller-origin ME (ME-S3, ME-S4) events are
instead tagged with a sentinel physical-sender identity meaning "fabricated
internally by the controller, no network packet involved," which trips a
bypass check that skips Stage~0 entirely — since the controller would
otherwise be verifying its own fabrication — so these two scenarios always
reach the TGN as attack-labelled events and alone supply the model's
entire ME-class training signal. This directly explains the ME/data-plane
cell's MCC=1.000$^\dagger$ and AUROC=1.000$^\dagger$ ($\epsilon$-smoothed,
footnoted) in Table~\ref{tab:tgn-val2-cross-mcc} (only 73 test rows, all
benign after Stage-0 filtering, so no positive class exists in that
specific subset to score against — 2 false positives were still recorded
among the 71 true negatives, Accuracy=0.973) and continues to be recorded
as a structural limitation of the pipeline's crypto layer correctly
rejecting implausible witness claims, rather than a
training defect.
